\section{Conclusion}
\label{sec:conclusion}

We systematically investigate category-completion false memories in LLMs, testing whether models falsely assert the existence of plausible-but-nonexistent items within well-populated categories.
Using the nonexistent seahorse emoji as a starting point, we probe two commercial LLMs across 162~items in 7~categories with 4~prompt framings.

Our findings paint a nuanced picture.
LLMs do exhibit category-completion false memories---plausible nonexistent items are affirmed at higher rates than implausible ones---but the driving mechanism is item-level semantic plausibility rather than category density.
The smallest category tested (insects) produces the highest error rates, directly contradicting the density hypothesis.
Adversarial framing amplifies false positive rates from 30\% to 80\%, demonstrating that conversational presuppositions can reliably override factual knowledge.
The phenomenon is domain-specific: emojis are highly vulnerable while structured factual knowledge remains robust.

These results have implications for LLM deployment.
Emoji-like knowledge domains---mutable catalogs where training data contains contradictory claims---require special attention in evaluation and post-training calibration.
Users and system designers should be aware that leading questions dramatically increase confabulation rates.
Future work should scale the probe dataset, extend to additional models, and apply mechanistic interpretability methods to understand why specific items trigger false memories while others do not.
